% =============================================================================
%  2-Desarrollo.tex — Bitácora de avances y autoevaluación
%
%  INSTRUCCIONES GENERALES:
%  ─────────────────────────────────────────────────────────────────────────────
%  Este archivo contiene dos tablas pre-configuradas:
%
%  TABLA 1 (longtable) — Bitácora cronológica de commits/tags
%    • Llena una fila por cada hito significativo del git log.
%    • Agrupa las filas por semana o etapa usando las filas de encabezado
%      de color (multicolumn). Añade o elimina semanas según necesites.
%    • Obtén los hashes y fechas con:
%        cd <TESIS_DIR> && git log --oneline --format="%h %ad %s" --date=short
%
%  TABLA 2 (tabularx) — Autoevaluación por capítulo o módulo
%    • Ajusta las filas a los capítulos/módulos de tu repositorio.
%    • Las columnas son: Capítulo | Avance concreto | Dificultad | Estado actual
%
%  Parámetros útiles disponibles:
%    \NombreRepo, \VerRef, \VerActual  (definidos en main.tex)
% =============================================================================

\section{Desarrollo y Avances}
\label{sec:desarrollo}

A través del sistema de control de versiones del repositorio \texttt{\NombreRepo},
es posible rastrear con precisión la evolución del proyecto entre las versiones
\texttt{\VerRef} y \texttt{\VerActual}. El Cuadro~\ref{tab:bitacora_git}
presenta una bitácora cronológica de los hitos registrados.

% ─────────────────────────────────────────────────────────────────────────────
%  TABLA 1: Bitácora cronológica (git log)
%
%  Anchos de columna: suma + 2×\tabcolsep×4cols debe ser ≤ \textwidth (155mm)
%  Con \tabcolsep=6pt: 4cols×12pt = 48pt = 16.9mm → columnas ≤ 138mm = 13.8cm
%  Distribución actual: 2.4+2.0+3.3+5.8 = 13.5cm → total 13.5+1.69 = 15.19cm ✓
% ─────────────────────────────────────────────────────────────────────────────
\begin{longtable}{>{\bfseries}p{2.4cm} p{2.0cm} p{3.3cm} p{5.8cm}}
    \caption{Bitácora cronológica de avances en \texttt{\NombreRepo}
             (Corte: \FechaEntrega).}
    \label{tab:bitacora_git} \\
    \toprule
    \textbf{Versión / Hito} & \textbf{Fecha} & \textbf{Ámbito} & \textbf{Descripción del avance} \\
    \midrule
    \endfirsthead

    \multicolumn{4}{l}{\footnotesize\textit{(Continuación del Cuadro~\ref{tab:bitacora_git})}} \\
    \toprule
    \textbf{Versión / Hito} & \textbf{Fecha} & \textbf{Ámbito} & \textbf{Descripción del avance} \\
    \midrule
    \endhead

    \midrule
    \multicolumn{4}{r}{\footnotesize\textit{Continúa en la siguiente página}} \\
    \endfoot

    \bottomrule
    \endlastfoot

    % ── Etapa 1 ───────────────────────────────────────────────────────────────
    % Cambia el texto de la etiqueta y las fechas según tu proyecto.
    \multicolumn{4}{l}{\small\color{ColorPrincipal}\textit{Etapa 1 — Descripción (dd--dd de mes de año)}} \\
    \addlinespace[2pt]

    % Fila de ejemplo — duplica o elimina según tus commits
    Hito o tag
    & DD mmm AAAA
    & Capítulo / Módulo
    & Descripción del cambio realizado (\texttt{abc1234}). \\
    \addlinespace[3pt]

    Hito o tag
    & DD mmm AAAA
    & Capítulo / Módulo
    & Descripción del cambio realizado (\texttt{def5678}). \\
    \addlinespace[8pt]

    % ── Etapa 2 ───────────────────────────────────────────────────────────────
    \multicolumn{4}{l}{\small\color{ColorPrincipal}\textit{Etapa 2 — Descripción (dd--dd de mes de año)}} \\
    \addlinespace[2pt]

    \texttt{v0.X.Y} — Nombre
    & DD mmm AAAA
    & Capítulo / Módulo
    & \textbf{Hito importante}: descripción detallada del avance o redefinición
      (\texttt{abc1234}). \\
    \addlinespace[8pt]

    % ── Etapa 3 ───────────────────────────────────────────────────────────────
    \multicolumn{4}{l}{\small\color{ColorPrincipal}\textit{Etapa 3 — Descripción (dd--dd de mes de año)}} \\
    \addlinespace[2pt]

    \texttt{v0.X.Z} — Nombre
    & DD mmm AAAA
    & Capítulo / Módulo
    & Descripción (\texttt{abc1234}). \\

\end{longtable}

% ─────────────────────────────────────────────────────────────────────────────
%  TABLA 2: Autoevaluación por capítulo / módulo
% ─────────────────────────────────────────────────────────────────────────────
\vspace{1em}

El Cuadro~\ref{tab:autoevaluacion} presenta una autoevaluación del estado actual
por capítulo o módulo, identificando avances, dificultades encontradas y
los apartados que requieren atención en la siguiente etapa.

\begin{table}[H]
    \centering
    \caption{Autoevaluación por capítulo: avances, redefiniciones y apartados débiles.}
    \label{tab:autoevaluacion}
    \fontsize{9pt}{11.5pt}\selectfont
    \begin{tabularx}{\textwidth}{
        >{\hsize=0.55\hsize\raggedright\arraybackslash}X
        >{\hsize=1.15\hsize\raggedright\arraybackslash}X
        >{\hsize=1.15\hsize\raggedright\arraybackslash}X
        >{\hsize=1.15\hsize\raggedright\arraybackslash}X}
        \toprule
        \textbf{Capítulo / Módulo} &
        \textbf{Avance concreto} &
        \textbf{Dificultad / Redefinición} &
        \textbf{Estado actual y pendientes} \\
        \midrule

        % Fila de ejemplo — duplica o elimina según tus módulos
        \textbf{Cap. 1}\newline Nombre del capítulo &
        Descripción del avance realizado. &
        Descripción de la dificultad o cambio de enfoque. &
        \textit{Estado:} descripción.\newline
        \textit{Pendiente:} descripción. \\
        \addlinespace\hline\addlinespace

        \textbf{Cap. 2}\newline Nombre del capítulo &
        Descripción del avance realizado. &
        Descripción de la dificultad. &
        \textit{Estado:} descripción.\newline
        \textit{Pendiente:} descripción. \\
        \bottomrule
    \end{tabularx}
\end{table}

\vspace{1em}
\textit{Nota metodológica:} Para respaldar esta autoevaluación, en la sección
siguiente se presenta evidencia generada con \texttt{latexdiff}, donde las
adiciones aparecen en \textcolor{blue}{azul} y las eliminaciones en
\textcolor{red}{rojo}.
